\section*{Lecture 1 - Overview}
\addcontentsline{toc}{section}{Lecture 1 - Overview}

\clearpage
\SlideHeader{Lecture 1 - Overview}{001}{表紙}
\begin{minipage}[t]{0.405\textwidth}
\SlideImage{1}{../Materials/2026/3DM_01_Overview_SS26.pdf}
\begin{adjustbox}{max totalsize={\textwidth}{0.43\textheight},center}
\begin{minipage}{\textwidth}\scriptsize
\NoteSection{日本語訳}\begin{itemize}
\item Data Driven Decision Making
\item 第1回 - 概要
\end{itemize}
\end{minipage}
\end{adjustbox}
\end{minipage}\hfill
\begin{minipage}[t]{0.58\textwidth}
\begin{adjustbox}{max totalsize={\textwidth}{0.89\textheight},center}
\begin{minipage}{\textwidth}\scriptsize
\NoteSection{注記}このページは表紙・目次・事務情報・導入画像が中心であるため、無理に確認問題や過去問を付けない。
\end{minipage}
\end{adjustbox}
\end{minipage}


\clearpage
\SlideHeader{Lecture 1 - Overview}{002}{連絡先}
\begin{minipage}[t]{0.405\textwidth}
\SlideImage{2}{../Materials/2026/3DM_01_Overview_SS26.pdf}
\begin{adjustbox}{max totalsize={\textwidth}{0.43\textheight},center}
\begin{minipage}{\textwidth}\scriptsize
\NoteSection{日本語訳}\begin{itemize}
\item 担当教員
\item ティーチングアシスタント
\item Decision Support オフィス
\item Webサイト
\end{itemize}
\end{minipage}
\end{adjustbox}
\end{minipage}\hfill
\begin{minipage}[t]{0.58\textwidth}
\begin{adjustbox}{max totalsize={\textwidth}{0.89\textheight},center}
\begin{minipage}{\textwidth}\scriptsize
\NoteSection{注記}このページは表紙・目次・事務情報・導入画像が中心であるため、無理に確認問題や過去問を付けない。
\end{minipage}
\end{adjustbox}
\end{minipage}


\clearpage
\SlideHeader{Lecture 1 - Overview}{003}{Decision Support の修士科目}
\begin{minipage}[t]{0.405\textwidth}
\SlideImage{3}{../Materials/2026/3DM_01_Overview_SS26.pdf}
\begin{adjustbox}{max totalsize={\textwidth}{0.43\textheight},center}
\begin{minipage}{\textwidth}\scriptsize
\NoteSection{日本語訳}\begin{itemize}
\item 基礎・方向づけ
\item 専門
\item その他の専攻、とくに機械工学系学生にも開かれている。
\end{itemize}
\end{minipage}
\end{adjustbox}
\end{minipage}\hfill
\begin{minipage}[t]{0.58\textwidth}
\begin{adjustbox}{max totalsize={\textwidth}{0.89\textheight},center}
\begin{minipage}{\textwidth}\scriptsize
\NoteSection{注記}このページは表紙・目次・事務情報・導入画像が中心であるため、無理に確認問題や過去問を付けない。
\end{minipage}
\end{adjustbox}
\end{minipage}


\clearpage
\SlideHeader{Lecture 1 - Overview}{004}{講義について}
\begin{minipage}[t]{0.405\textwidth}
\SlideImage{4}{../Materials/2026/3DM_01_Overview_SS26.pdf}
\begin{adjustbox}{max totalsize={\textwidth}{0.43\textheight},center}
\begin{minipage}{\textwidth}\scriptsize
\NoteSection{日本語訳}\begin{itemize}
\item Data Driven Decision Making
\item 以前は『モビリティ応用のための情報システム』、さらに以前は『交通情報システム』だった。
\item 焦点は、情報システムの描写から、情報システム内の自動意思決定へ移っている。
\item Uber, Lime, Amazon などの先端的デジタルビジネスモデルと整合する。
\item PMT は意思決定に、IDA はデータ分析に焦点を置く。DDDM はその上に構築される。
\item 3DM は演習とともに『Decision Support専門』を構成する。
\end{itemize}\NoteSection{試験対策ポイント}\begin{itemize}
\item DDDMはデータ分析を意思決定と実装へ接続する講義である。
\item stochasticity と dynamism を分けて説明できることが最重要である。
\end{itemize}
\end{minipage}
\end{adjustbox}
\end{minipage}\hfill
\begin{minipage}[t]{0.58\textwidth}
\begin{adjustbox}{max totalsize={\textwidth}{0.89\textheight},center}
\begin{minipage}{\textwidth}\scriptsize
\NoteSection{解説}このページでは「講義について」を理解する。スライド上の表現は短いが、試験では定義、具体例、モデル上の意味、手法の限界まで説明できる必要がある。\par
Data Driven Decision Making は、単にデータを分析してレポートを作ることではない。データを使って、現実で実行する行動を選び、その結果をまた観測して次の意思決定へ反映する考え方である。たとえば食品配送であれば、注文数を予測するだけではなく、どの注文を受けるか、どのドライバーに割り当てるか、どの順番で回るかまで決める。\par
この講義の中心は、現実の意思決定が静的で完全情報の問題ではないという点にある。現実では、注文、交通、ドライバー位置、キャンセル、サービス時間などが時間とともに変わる。したがって、最初に作った計画を最後まで固定するのではなく、新しい情報を受けて状態を更新し、必要に応じて計画を変える必要がある。\par
stochasticity は、将来何が起こるかが確率的・不確実であることを意味する。dynamism は、時間が進むにつれて状態が変わり、意思決定を繰り返す必要があることを意味する。試験では、この二語を同義語のように書かないことが重要である。注文がいつ来るか分からないことは stochasticity、注文が来た後に計画を更新することは dynamism である。\par
答案では、DDDMを Business Analytics、Operations Research、Data Analytics、情報システム実装の交点として説明すると強い。予測モデルだけでは『何をするか』が決まらず、最適化モデルだけでは『現実データをどう入れるか』が決まらない。DDDMはこの間を接続し、現実のデータから実行可能な意思決定を作る。\par
具体例として、Uberや食品配送では、リアルタイム注文、ドライバーの現在位置、交通状況、顧客の待ち時間を使いながら、数分ごとに新しい割当を作る。このとき、過去データから需要を学習するだけでなく、将来の不確実な注文を見越して現在のドライバーを使い切らないようにするなど、将来の柔軟性も重要になる。\par\NoteSection{確認問題と解答}\begin{enumerate}\item \textbf{L01-S004: 「講義について」の概念を、定義・具体例・モデル上の意味の順に説明せよ。}\\定義としては、DDDMはデータ分析を意思決定と実装へ接続する講義である。。具体例では、配送・在庫・フードデリバリーのような現実問題を用い、状態、決定、外生情報、報酬、または情報モデル/意思決定モデルのどこに対応するかを明示する。モデル上の意味として、この概念が目的関数、制約、状態表現、将来価値、または方策選択にどのような影響を与えるかを書く。試験では、用語だけを列挙せず、入力データから意思決定、さらに現実の実装へつながる流れを文章で示すことが重要である。\item \textbf{L01-S004: このスライドの考え方を、講義全体のDDDMプロセスの中に位置づけよ。}\\DDDMプロセスでは、現実世界のデータを集約して情報モデルを作り、意思決定モデルまたは方策で行動を選び、実装後に結果を観測して次の状態へ進む。このスライドの概念は、そのどこか一部だけではなく、データ表現、意思決定、将来不確実性、計算可能性のいずれかに関わる。答案では、まずこの概念がどの段階に属するかを述べ、次に他の段階へどう影響するかを書く。例えば価値関数なら将来価値を現在決定へ戻す役割、FIFOなら旅行時間情報モデルの整合性を保証する役割、CFAなら目的関数を調整して将来に良い行動を誘導する役割である。\end{enumerate}
\end{minipage}
\end{adjustbox}
\end{minipage}


\clearpage
\SlideHeader{Lecture 1 - Overview}{005}{目標}
\begin{minipage}[t]{0.405\textwidth}
\SlideImage{5}{../Materials/2026/3DM_01_Overview_SS26.pdf}
\begin{adjustbox}{max totalsize={\textwidth}{0.43\textheight},center}
\begin{minipage}{\textwidth}\scriptsize
\NoteSection{日本語訳}\begin{itemize}
\item 確率性と動的性質を理解する。
\item 確率的・動的意思決定問題をモデル化する。
\item 方法論の概要、機能、強み、弱みを理解する。
\item デジタル経済における関連問題を理解する。
\item 60分筆記試験に合格する。
\item 修士論文、博士課程への準備を行う。
\end{itemize}\NoteSection{試験対策ポイント}\begin{itemize}
\item DDDMはデータ分析を意思決定と実装へ接続する講義である。
\item stochasticity と dynamism を分けて説明できることが最重要である。
\end{itemize}
\end{minipage}
\end{adjustbox}
\end{minipage}\hfill
\begin{minipage}[t]{0.58\textwidth}
\begin{adjustbox}{max totalsize={\textwidth}{0.89\textheight},center}
\begin{minipage}{\textwidth}\scriptsize
\NoteSection{解説}このページでは「目標」を理解する。スライド上の表現は短いが、試験では定義、具体例、モデル上の意味、手法の限界まで説明できる必要がある。\par
Data Driven Decision Making は、単にデータを分析してレポートを作ることではない。データを使って、現実で実行する行動を選び、その結果をまた観測して次の意思決定へ反映する考え方である。たとえば食品配送であれば、注文数を予測するだけではなく、どの注文を受けるか、どのドライバーに割り当てるか、どの順番で回るかまで決める。\par
この講義の中心は、現実の意思決定が静的で完全情報の問題ではないという点にある。現実では、注文、交通、ドライバー位置、キャンセル、サービス時間などが時間とともに変わる。したがって、最初に作った計画を最後まで固定するのではなく、新しい情報を受けて状態を更新し、必要に応じて計画を変える必要がある。\par
stochasticity は、将来何が起こるかが確率的・不確実であることを意味する。dynamism は、時間が進むにつれて状態が変わり、意思決定を繰り返す必要があることを意味する。試験では、この二語を同義語のように書かないことが重要である。注文がいつ来るか分からないことは stochasticity、注文が来た後に計画を更新することは dynamism である。\par
答案では、DDDMを Business Analytics、Operations Research、Data Analytics、情報システム実装の交点として説明すると強い。予測モデルだけでは『何をするか』が決まらず、最適化モデルだけでは『現実データをどう入れるか』が決まらない。DDDMはこの間を接続し、現実のデータから実行可能な意思決定を作る。\par
具体例として、Uberや食品配送では、リアルタイム注文、ドライバーの現在位置、交通状況、顧客の待ち時間を使いながら、数分ごとに新しい割当を作る。このとき、過去データから需要を学習するだけでなく、将来の不確実な注文を見越して現在のドライバーを使い切らないようにするなど、将来の柔軟性も重要になる。\par\NoteSection{確認問題と解答}\begin{enumerate}\item \textbf{L01-S005: 「目標」の概念を、定義・具体例・モデル上の意味の順に説明せよ。}\\定義としては、DDDMはデータ分析を意思決定と実装へ接続する講義である。。具体例では、配送・在庫・フードデリバリーのような現実問題を用い、状態、決定、外生情報、報酬、または情報モデル/意思決定モデルのどこに対応するかを明示する。モデル上の意味として、この概念が目的関数、制約、状態表現、将来価値、または方策選択にどのような影響を与えるかを書く。試験では、用語だけを列挙せず、入力データから意思決定、さらに現実の実装へつながる流れを文章で示すことが重要である。\item \textbf{L01-S005: このスライドの考え方を、講義全体のDDDMプロセスの中に位置づけよ。}\\DDDMプロセスでは、現実世界のデータを集約して情報モデルを作り、意思決定モデルまたは方策で行動を選び、実装後に結果を観測して次の状態へ進む。このスライドの概念は、そのどこか一部だけではなく、データ表現、意思決定、将来不確実性、計算可能性のいずれかに関わる。答案では、まずこの概念がどの段階に属するかを述べ、次に他の段階へどう影響するかを書く。例えば価値関数なら将来価値を現在決定へ戻す役割、FIFOなら旅行時間情報モデルの整合性を保証する役割、CFAなら目的関数を調整して将来に良い行動を誘導する役割である。\end{enumerate}
\end{minipage}
\end{adjustbox}
\end{minipage}


\clearpage
\SlideHeader{Lecture 1 - Overview}{006}{DDDMの講義計画}
\begin{minipage}[t]{0.405\textwidth}
\SlideImage{6}{../Materials/2026/3DM_01_Overview_SS26.pdf}
\begin{adjustbox}{max totalsize={\textwidth}{0.43\textheight},center}
\begin{minipage}{\textwidth}\scriptsize
\NoteSection{日本語訳}\begin{itemize}
\item 講義トピック
\end{itemize}
\end{minipage}
\end{adjustbox}
\end{minipage}\hfill
\begin{minipage}[t]{0.58\textwidth}
\begin{adjustbox}{max totalsize={\textwidth}{0.89\textheight},center}
\begin{minipage}{\textwidth}\scriptsize
\NoteSection{注記}このページは表紙・目次・事務情報・導入画像が中心であるため、無理に確認問題や過去問を付けない。
\end{minipage}
\end{adjustbox}
\end{minipage}


\clearpage
\SlideHeader{Lecture 1 - Overview}{007}{文献}
\begin{minipage}[t]{0.405\textwidth}
\SlideImage{7}{../Materials/2026/3DM_01_Overview_SS26.pdf}
\begin{adjustbox}{max totalsize={\textwidth}{0.43\textheight},center}
\begin{minipage}{\textwidth}\scriptsize
\NoteSection{日本語訳}\begin{itemize}
\item Warren Powell の研究。
\item TU Braunschweig における Marlin Ulmer らとの近年の共同研究。
\item 交通・モビリティに関する適用と発展。
\item Ulmer et al. (2020): stochastic dynamic vehicle routing problems のモデル化。
\item Ulmer et al. (2021): prescriptive analytics の観点から見た stochastic dynamic vehicle routing のレビュー。
\end{itemize}
\end{minipage}
\end{adjustbox}
\end{minipage}\hfill
\begin{minipage}[t]{0.58\textwidth}
\begin{adjustbox}{max totalsize={\textwidth}{0.89\textheight},center}
\begin{minipage}{\textwidth}\scriptsize
\NoteSection{注記}このページは表紙・目次・事務情報・導入画像が中心であるため、無理に確認問題や過去問を付けない。
\end{minipage}
\end{adjustbox}
\end{minipage}


\clearpage
\SlideHeader{Lecture 1 - Overview}{008}{演習について}
\begin{minipage}[t]{0.405\textwidth}
\SlideImage{8}{../Materials/2026/3DM_01_Overview_SS26.pdf}
\begin{adjustbox}{max totalsize={\textwidth}{0.43\textheight},center}
\begin{minipage}{\textwidth}\scriptsize
\NoteSection{日本語訳}\begin{itemize}
\item 『Decision Support専門』の Studienleistung (2.5 LP)。
\item serious game アプリケーション『Trucks and Barges』で講義内容を練習する。
\item シミュレーション内容
\item 運営
\end{itemize}
\end{minipage}
\end{adjustbox}
\end{minipage}\hfill
\begin{minipage}[t]{0.58\textwidth}
\begin{adjustbox}{max totalsize={\textwidth}{0.89\textheight},center}
\begin{minipage}{\textwidth}\scriptsize
\NoteSection{注記}このページは表紙・目次・事務情報・導入画像が中心であるため、無理に確認問題や過去問を付けない。
\end{minipage}
\end{adjustbox}
\end{minipage}


\clearpage
\SlideHeader{Lecture 1 - Overview}{009}{試験情報}
\begin{minipage}[t]{0.405\textwidth}
\SlideImage{9}{../Materials/2026/3DM_01_Overview_SS26.pdf}
\begin{adjustbox}{max totalsize={\textwidth}{0.43\textheight},center}
\begin{minipage}{\textwidth}\scriptsize
\NoteSection{日本語訳}\begin{itemize}
\item 試験会場と時刻は講座のWebサイトで確認できる。
\item 情報は定期的に更新される。
\item 試験日は早めに告知され、試験教室は試験直前に割り当てられる。
\item 必要な情報がWebサイトにない場合のみメールで問い合わせる。
\end{itemize}
\end{minipage}
\end{adjustbox}
\end{minipage}\hfill
\begin{minipage}[t]{0.58\textwidth}
\begin{adjustbox}{max totalsize={\textwidth}{0.89\textheight},center}
\begin{minipage}{\textwidth}\scriptsize
\NoteSection{注記}このページは表紙・目次・事務情報・導入画像が中心であるため、無理に確認問題や過去問を付けない。
\end{minipage}
\end{adjustbox}
\end{minipage}


\clearpage
\SlideHeader{Lecture 1 - Overview}{010}{過去問}
\begin{minipage}[t]{0.405\textwidth}
\SlideImage{10}{../Materials/2026/3DM_01_Overview_SS26.pdf}
\begin{adjustbox}{max totalsize={\textwidth}{0.43\textheight},center}
\begin{minipage}{\textwidth}\scriptsize
\NoteSection{日本語訳}\begin{itemize}
\item 近年の過去問の印刷版は研究所で入手できる。部数には限りがある。
\item 過去問はWebサイトでも確認できる。
\item 過去問を解くだけでは十分な準備ではなく、良い成績を保証しない。
\item 講義で扱った全トピックが試験対象である。
\item 問題形式は似ることがあるが、内容は変わり、過去問にない新問も含まれ得る。
\item 過去問の模範解答は提供されない。
\end{itemize}
\end{minipage}
\end{adjustbox}
\end{minipage}\hfill
\begin{minipage}[t]{0.58\textwidth}
\begin{adjustbox}{max totalsize={\textwidth}{0.89\textheight},center}
\begin{minipage}{\textwidth}\scriptsize
\NoteSection{注記}このページは表紙・目次・事務情報・導入画像が中心であるため、無理に確認問題や過去問を付けない。
\end{minipage}
\end{adjustbox}
\end{minipage}


\clearpage
\SlideHeader{Lecture 1 - Overview}{011}{アジェンダ}
\begin{minipage}[t]{0.405\textwidth}
\SlideImage{11}{../Materials/2026/3DM_01_Overview_SS26.pdf}
\begin{adjustbox}{max totalsize={\textwidth}{0.43\textheight},center}
\begin{minipage}{\textwidth}\scriptsize
\NoteSection{日本語訳}\begin{itemize}
\item 動機
\item Operations Research と Business Analytics の復習
\item モデリングと講義全体の見取り図
\end{itemize}
\end{minipage}
\end{adjustbox}
\end{minipage}\hfill
\begin{minipage}[t]{0.58\textwidth}
\begin{adjustbox}{max totalsize={\textwidth}{0.89\textheight},center}
\begin{minipage}{\textwidth}\scriptsize
\NoteSection{注記}このページは表紙・目次・事務情報・導入画像が中心であるため、無理に確認問題や過去問を付けない。
\end{minipage}
\end{adjustbox}
\end{minipage}


\clearpage
\SlideHeader{Lecture 1 - Overview}{012}{動機}
\begin{minipage}[t]{0.405\textwidth}
\SlideImage{12}{../Materials/2026/3DM_01_Overview_SS26.pdf}
\begin{adjustbox}{max totalsize={\textwidth}{0.43\textheight},center}
\begin{minipage}{\textwidth}\scriptsize
\NoteSection{日本語訳}\begin{itemize}
\item 新しく変化するビジネスコンセプト: より速く、より個別的に、リアルタイムに。
\item 例
\item 計画と制御が難しくなる。
\item 二つの共通テーマ
\item 確率的・動的意思決定過程へつながる。
\end{itemize}\NoteSection{試験対策ポイント}\begin{itemize}
\item DDDMはデータ分析を意思決定と実装へ接続する講義である。
\item stochasticity と dynamism を分けて説明できることが最重要である。
\end{itemize}
\end{minipage}
\end{adjustbox}
\end{minipage}\hfill
\begin{minipage}[t]{0.58\textwidth}
\begin{adjustbox}{max totalsize={\textwidth}{0.89\textheight},center}
\begin{minipage}{\textwidth}\scriptsize
\NoteSection{解説}このページでは「動機」を理解する。スライド上の表現は短いが、試験では定義、具体例、モデル上の意味、手法の限界まで説明できる必要がある。\par
Data Driven Decision Making は、単にデータを分析してレポートを作ることではない。データを使って、現実で実行する行動を選び、その結果をまた観測して次の意思決定へ反映する考え方である。たとえば食品配送であれば、注文数を予測するだけではなく、どの注文を受けるか、どのドライバーに割り当てるか、どの順番で回るかまで決める。\par
この講義の中心は、現実の意思決定が静的で完全情報の問題ではないという点にある。現実では、注文、交通、ドライバー位置、キャンセル、サービス時間などが時間とともに変わる。したがって、最初に作った計画を最後まで固定するのではなく、新しい情報を受けて状態を更新し、必要に応じて計画を変える必要がある。\par
stochasticity は、将来何が起こるかが確率的・不確実であることを意味する。dynamism は、時間が進むにつれて状態が変わり、意思決定を繰り返す必要があることを意味する。試験では、この二語を同義語のように書かないことが重要である。注文がいつ来るか分からないことは stochasticity、注文が来た後に計画を更新することは dynamism である。\par
答案では、DDDMを Business Analytics、Operations Research、Data Analytics、情報システム実装の交点として説明すると強い。予測モデルだけでは『何をするか』が決まらず、最適化モデルだけでは『現実データをどう入れるか』が決まらない。DDDMはこの間を接続し、現実のデータから実行可能な意思決定を作る。\par
具体例として、Uberや食品配送では、リアルタイム注文、ドライバーの現在位置、交通状況、顧客の待ち時間を使いながら、数分ごとに新しい割当を作る。このとき、過去データから需要を学習するだけでなく、将来の不確実な注文を見越して現在のドライバーを使い切らないようにするなど、将来の柔軟性も重要になる。\par\NoteSection{確認問題と解答}\begin{enumerate}\item \textbf{L01-S012: 「動機」の概念を、定義・具体例・モデル上の意味の順に説明せよ。}\\定義としては、DDDMはデータ分析を意思決定と実装へ接続する講義である。。具体例では、配送・在庫・フードデリバリーのような現実問題を用い、状態、決定、外生情報、報酬、または情報モデル/意思決定モデルのどこに対応するかを明示する。モデル上の意味として、この概念が目的関数、制約、状態表現、将来価値、または方策選択にどのような影響を与えるかを書く。試験では、用語だけを列挙せず、入力データから意思決定、さらに現実の実装へつながる流れを文章で示すことが重要である。\item \textbf{L01-S012: このスライドの考え方を、講義全体のDDDMプロセスの中に位置づけよ。}\\DDDMプロセスでは、現実世界のデータを集約して情報モデルを作り、意思決定モデルまたは方策で行動を選び、実装後に結果を観測して次の状態へ進む。このスライドの概念は、そのどこか一部だけではなく、データ表現、意思決定、将来不確実性、計算可能性のいずれかに関わる。答案では、まずこの概念がどの段階に属するかを述べ、次に他の段階へどう影響するかを書く。例えば価値関数なら将来価値を現在決定へ戻す役割、FIFOなら旅行時間情報モデルの整合性を保証する役割、CFAなら目的関数を調整して将来に良い行動を誘導する役割である。\end{enumerate}
\end{minipage}
\end{adjustbox}
\end{minipage}


\clearpage
\SlideHeader{Lecture 1 - Overview}{013}{例: 輸送}
\begin{minipage}[t]{0.405\textwidth}
\SlideImage{13}{../Materials/2026/3DM_01_Overview_SS26.pdf}
\begin{adjustbox}{max totalsize={\textwidth}{0.43\textheight},center}
\begin{minipage}{\textwidth}\scriptsize
\NoteSection{日本語訳}\begin{itemize}
\item 同日配送、ピックアップ、ドローン配送、食品配送など、輸送サービスの例。
\end{itemize}\NoteSection{試験対策ポイント}\begin{itemize}
\item DDDMはデータ分析を意思決定と実装へ接続する講義である。
\item stochasticity と dynamism を分けて説明できることが最重要である。
\end{itemize}
\end{minipage}
\end{adjustbox}
\end{minipage}\hfill
\begin{minipage}[t]{0.58\textwidth}
\begin{adjustbox}{max totalsize={\textwidth}{0.89\textheight},center}
\begin{minipage}{\textwidth}\scriptsize
\NoteSection{解説}このページでは「例: 輸送」を理解する。スライド上の表現は短いが、試験では定義、具体例、モデル上の意味、手法の限界まで説明できる必要がある。\par
Data Driven Decision Making は、単にデータを分析してレポートを作ることではない。データを使って、現実で実行する行動を選び、その結果をまた観測して次の意思決定へ反映する考え方である。たとえば食品配送であれば、注文数を予測するだけではなく、どの注文を受けるか、どのドライバーに割り当てるか、どの順番で回るかまで決める。\par
この講義の中心は、現実の意思決定が静的で完全情報の問題ではないという点にある。現実では、注文、交通、ドライバー位置、キャンセル、サービス時間などが時間とともに変わる。したがって、最初に作った計画を最後まで固定するのではなく、新しい情報を受けて状態を更新し、必要に応じて計画を変える必要がある。\par
stochasticity は、将来何が起こるかが確率的・不確実であることを意味する。dynamism は、時間が進むにつれて状態が変わり、意思決定を繰り返す必要があることを意味する。試験では、この二語を同義語のように書かないことが重要である。注文がいつ来るか分からないことは stochasticity、注文が来た後に計画を更新することは dynamism である。\par
答案では、DDDMを Business Analytics、Operations Research、Data Analytics、情報システム実装の交点として説明すると強い。予測モデルだけでは『何をするか』が決まらず、最適化モデルだけでは『現実データをどう入れるか』が決まらない。DDDMはこの間を接続し、現実のデータから実行可能な意思決定を作る。\par
具体例として、Uberや食品配送では、リアルタイム注文、ドライバーの現在位置、交通状況、顧客の待ち時間を使いながら、数分ごとに新しい割当を作る。このとき、過去データから需要を学習するだけでなく、将来の不確実な注文を見越して現在のドライバーを使い切らないようにするなど、将来の柔軟性も重要になる。\par\NoteSection{確認問題と解答}\begin{enumerate}\item \textbf{L01-S013: 「例: 輸送」の概念を、定義・具体例・モデル上の意味の順に説明せよ。}\\定義としては、DDDMはデータ分析を意思決定と実装へ接続する講義である。。具体例では、配送・在庫・フードデリバリーのような現実問題を用い、状態、決定、外生情報、報酬、または情報モデル/意思決定モデルのどこに対応するかを明示する。モデル上の意味として、この概念が目的関数、制約、状態表現、将来価値、または方策選択にどのような影響を与えるかを書く。試験では、用語だけを列挙せず、入力データから意思決定、さらに現実の実装へつながる流れを文章で示すことが重要である。\item \textbf{L01-S013: このスライドの考え方を、講義全体のDDDMプロセスの中に位置づけよ。}\\DDDMプロセスでは、現実世界のデータを集約して情報モデルを作り、意思決定モデルまたは方策で行動を選び、実装後に結果を観測して次の状態へ進む。このスライドの概念は、そのどこか一部だけではなく、データ表現、意思決定、将来不確実性、計算可能性のいずれかに関わる。答案では、まずこの概念がどの段階に属するかを述べ、次に他の段階へどう影響するかを書く。例えば価値関数なら将来価値を現在決定へ戻す役割、FIFOなら旅行時間情報モデルの整合性を保証する役割、CFAなら目的関数を調整して将来に良い行動を誘導する役割である。\end{enumerate}
\end{minipage}
\end{adjustbox}
\end{minipage}


\clearpage
\SlideHeader{Lecture 1 - Overview}{014}{例: モビリティとサービス}
\begin{minipage}[t]{0.405\textwidth}
\SlideImage{14}{../Materials/2026/3DM_01_Overview_SS26.pdf}
\begin{adjustbox}{max totalsize={\textwidth}{0.43\textheight},center}
\begin{minipage}{\textwidth}\scriptsize
\NoteSection{日本語訳}\begin{itemize}
\item バイクシェア、オンデマンド交通、医療訪問、タクシー、エレベーター保守、緊急車両など、モビリティ・サービスの例。
\end{itemize}\NoteSection{試験対策ポイント}\begin{itemize}
\item DDDMはデータ分析を意思決定と実装へ接続する講義である。
\item stochasticity と dynamism を分けて説明できることが最重要である。
\end{itemize}
\end{minipage}
\end{adjustbox}
\end{minipage}\hfill
\begin{minipage}[t]{0.58\textwidth}
\begin{adjustbox}{max totalsize={\textwidth}{0.89\textheight},center}
\begin{minipage}{\textwidth}\scriptsize
\NoteSection{解説}このページでは「例: モビリティとサービス」を理解する。スライド上の表現は短いが、試験では定義、具体例、モデル上の意味、手法の限界まで説明できる必要がある。\par
Data Driven Decision Making は、単にデータを分析してレポートを作ることではない。データを使って、現実で実行する行動を選び、その結果をまた観測して次の意思決定へ反映する考え方である。たとえば食品配送であれば、注文数を予測するだけではなく、どの注文を受けるか、どのドライバーに割り当てるか、どの順番で回るかまで決める。\par
この講義の中心は、現実の意思決定が静的で完全情報の問題ではないという点にある。現実では、注文、交通、ドライバー位置、キャンセル、サービス時間などが時間とともに変わる。したがって、最初に作った計画を最後まで固定するのではなく、新しい情報を受けて状態を更新し、必要に応じて計画を変える必要がある。\par
stochasticity は、将来何が起こるかが確率的・不確実であることを意味する。dynamism は、時間が進むにつれて状態が変わり、意思決定を繰り返す必要があることを意味する。試験では、この二語を同義語のように書かないことが重要である。注文がいつ来るか分からないことは stochasticity、注文が来た後に計画を更新することは dynamism である。\par
答案では、DDDMを Business Analytics、Operations Research、Data Analytics、情報システム実装の交点として説明すると強い。予測モデルだけでは『何をするか』が決まらず、最適化モデルだけでは『現実データをどう入れるか』が決まらない。DDDMはこの間を接続し、現実のデータから実行可能な意思決定を作る。\par
具体例として、Uberや食品配送では、リアルタイム注文、ドライバーの現在位置、交通状況、顧客の待ち時間を使いながら、数分ごとに新しい割当を作る。このとき、過去データから需要を学習するだけでなく、将来の不確実な注文を見越して現在のドライバーを使い切らないようにするなど、将来の柔軟性も重要になる。\par\NoteSection{確認問題と解答}\begin{enumerate}\item \textbf{L01-S014: 「例: モビリティとサービス」の概念を、定義・具体例・モデル上の意味の順に説明せよ。}\\定義としては、DDDMはデータ分析を意思決定と実装へ接続する講義である。。具体例では、配送・在庫・フードデリバリーのような現実問題を用い、状態、決定、外生情報、報酬、または情報モデル/意思決定モデルのどこに対応するかを明示する。モデル上の意味として、この概念が目的関数、制約、状態表現、将来価値、または方策選択にどのような影響を与えるかを書く。試験では、用語だけを列挙せず、入力データから意思決定、さらに現実の実装へつながる流れを文章で示すことが重要である。\item \textbf{L01-S014: このスライドの考え方を、講義全体のDDDMプロセスの中に位置づけよ。}\\DDDMプロセスでは、現実世界のデータを集約して情報モデルを作り、意思決定モデルまたは方策で行動を選び、実装後に結果を観測して次の状態へ進む。このスライドの概念は、そのどこか一部だけではなく、データ表現、意思決定、将来不確実性、計算可能性のいずれかに関わる。答案では、まずこの概念がどの段階に属するかを述べ、次に他の段階へどう影響するかを書く。例えば価値関数なら将来価値を現在決定へ戻す役割、FIFOなら旅行時間情報モデルの整合性を保証する役割、CFAなら目的関数を調整して将来に良い行動を誘導する役割である。\end{enumerate}
\end{minipage}
\end{adjustbox}
\end{minipage}


\clearpage
\SlideHeader{Lecture 1 - Overview}{015}{まとめ}
\begin{minipage}[t]{0.405\textwidth}
\SlideImage{15}{../Materials/2026/3DM_01_Overview_SS26.pdf}
\begin{adjustbox}{max totalsize={\textwidth}{0.43\textheight},center}
\begin{minipage}{\textwidth}\scriptsize
\NoteSection{日本語訳}\begin{itemize}
\item サービスプロセスは時間を通じて進行する。
\item 需要、依頼、交通、ドライバー、バッテリー残量、要件などの新情報が発生する。
\item 計画への反応と更新が必要になる。
\item 反応が遅すぎる場合がある。例: 車両が渋滞に巻き込まれる、食事配送のドライバーが周辺にいない、ドローンが墜落する。
\item 期待される将来を先読みする必要がある。
\item 将来のシステム状態に対する柔軟性を目指す。
\end{itemize}\NoteSection{試験対策ポイント}\begin{itemize}
\item DDDMはデータ分析を意思決定と実装へ接続する講義である。
\item stochasticity と dynamism を分けて説明できることが最重要である。
\end{itemize}
\end{minipage}
\end{adjustbox}
\end{minipage}\hfill
\begin{minipage}[t]{0.58\textwidth}
\begin{adjustbox}{max totalsize={\textwidth}{0.89\textheight},center}
\begin{minipage}{\textwidth}\scriptsize
\NoteSection{解説}このページでは「まとめ」を理解する。スライド上の表現は短いが、試験では定義、具体例、モデル上の意味、手法の限界まで説明できる必要がある。\par
Data Driven Decision Making は、単にデータを分析してレポートを作ることではない。データを使って、現実で実行する行動を選び、その結果をまた観測して次の意思決定へ反映する考え方である。たとえば食品配送であれば、注文数を予測するだけではなく、どの注文を受けるか、どのドライバーに割り当てるか、どの順番で回るかまで決める。\par
この講義の中心は、現実の意思決定が静的で完全情報の問題ではないという点にある。現実では、注文、交通、ドライバー位置、キャンセル、サービス時間などが時間とともに変わる。したがって、最初に作った計画を最後まで固定するのではなく、新しい情報を受けて状態を更新し、必要に応じて計画を変える必要がある。\par
stochasticity は、将来何が起こるかが確率的・不確実であることを意味する。dynamism は、時間が進むにつれて状態が変わり、意思決定を繰り返す必要があることを意味する。試験では、この二語を同義語のように書かないことが重要である。注文がいつ来るか分からないことは stochasticity、注文が来た後に計画を更新することは dynamism である。\par
答案では、DDDMを Business Analytics、Operations Research、Data Analytics、情報システム実装の交点として説明すると強い。予測モデルだけでは『何をするか』が決まらず、最適化モデルだけでは『現実データをどう入れるか』が決まらない。DDDMはこの間を接続し、現実のデータから実行可能な意思決定を作る。\par
具体例として、Uberや食品配送では、リアルタイム注文、ドライバーの現在位置、交通状況、顧客の待ち時間を使いながら、数分ごとに新しい割当を作る。このとき、過去データから需要を学習するだけでなく、将来の不確実な注文を見越して現在のドライバーを使い切らないようにするなど、将来の柔軟性も重要になる。\par\NoteSection{確認問題と解答}\begin{enumerate}\item \textbf{L01-S015: 「まとめ」の概念を、定義・具体例・モデル上の意味の順に説明せよ。}\\定義としては、DDDMはデータ分析を意思決定と実装へ接続する講義である。。具体例では、配送・在庫・フードデリバリーのような現実問題を用い、状態、決定、外生情報、報酬、または情報モデル/意思決定モデルのどこに対応するかを明示する。モデル上の意味として、この概念が目的関数、制約、状態表現、将来価値、または方策選択にどのような影響を与えるかを書く。試験では、用語だけを列挙せず、入力データから意思決定、さらに現実の実装へつながる流れを文章で示すことが重要である。\item \textbf{L01-S015: このスライドの考え方を、講義全体のDDDMプロセスの中に位置づけよ。}\\DDDMプロセスでは、現実世界のデータを集約して情報モデルを作り、意思決定モデルまたは方策で行動を選び、実装後に結果を観測して次の状態へ進む。このスライドの概念は、そのどこか一部だけではなく、データ表現、意思決定、将来不確実性、計算可能性のいずれかに関わる。答案では、まずこの概念がどの段階に属するかを述べ、次に他の段階へどう影響するかを書く。例えば価値関数なら将来価値を現在決定へ戻す役割、FIFOなら旅行時間情報モデルの整合性を保証する役割、CFAなら目的関数を調整して将来に良い行動を誘導する役割である。\end{enumerate}
\end{minipage}
\end{adjustbox}
\end{minipage}


\clearpage
\SlideHeader{Lecture 1 - Overview}{016}{意思決定}
\begin{minipage}[t]{0.405\textwidth}
\SlideImage{16}{../Materials/2026/3DM_01_Overview_SS26.pdf}
\begin{adjustbox}{max totalsize={\textwidth}{0.43\textheight},center}
\begin{minipage}{\textwidth}\scriptsize
\NoteSection{日本語訳}\begin{itemize}
\item Real-World: 行動が実行され、その結果として状態を観測できる。
\item Aggregation: 現実世界のデータを情報モデルの次元へ移す。
\item Information Model: 情報を圧縮された形で描写する。
\item Decision Model: 決定空間を定義し、決定を評価する。
\item Implementation: 意思決定モデルの解を現実世界の行動へ変換する。
\end{itemize}\NoteSection{試験対策ポイント}\begin{itemize}
\item Real-World, Aggregation, Information Model, Decision Model, Implementation の向きを正確に説明する。
\item 情報モデルは現実の圧縮表現、意思決定モデルは行動選択の評価構造である。
\end{itemize}
\end{minipage}
\end{adjustbox}
\end{minipage}\hfill
\begin{minipage}[t]{0.58\textwidth}
\begin{adjustbox}{max totalsize={\textwidth}{0.89\textheight},center}
\begin{minipage}{\textwidth}\scriptsize
\NoteSection{解説}このページでは「意思決定」を理解する。スライド上の表現は短いが、試験では定義、具体例、モデル上の意味、手法の限界まで説明できる必要がある。\par
Real-World は、配送車、顧客、道路、倉庫、ドライバー、注文、交通といった現実そのものである。現実は複雑すぎるため、そのまま数理モデルやアルゴリズムに入れることはできない。そこで、意思決定に必要な要素だけを抽出し、モデルが扱える形に変換する必要がある。\par
Aggregation は、現実から得られた詳細データを情報モデルの次元に圧縮する操作である。例として、GPS点列を道路セグメントの速度に変換する、住所を顧客ノードに変換する、過去走行ログから顧客間旅行時間行列を作る、注文履歴から需要分布を推定する、といった処理がある。\par
Information Model は、意思決定に必要な情報の構造化表現である。ここには、顧客集合、デポ、車両集合、旅行時間行列、需要、時間窓、確率分布などが含まれる。重要なのは、情報モデルは現実の全コピーではなく、意思決定に有用な抽象化だという点である。\par
Decision Model は、その情報モデルを入力として、どの行動を選べるか、何を最適化するか、どの制約を守るかを定義する。TSPなら、顧客を一度ずつ訪問し総旅行時間を最小化する。VRPなら、複数車両、容量制約、時間窓などが加わる。ここでは決定変数、目的関数、制約が中心になる。\par
Implementation は、モデル上の解を現実の行動に戻す処理である。たとえば x\_\{ij\}=1 というアーク選択を、ドライバーに提示する訪問順、住所、ナビゲーション、積み込み順へ変換する。実装後には実際の遅延や走行時間が観測され、それが再びAggregationを通じてモデル改善に使われる。\par\NoteSection{確認問題と解答}\begin{enumerate}\item \textbf{L01-S016: 「意思決定」の概念を、定義・具体例・モデル上の意味の順に説明せよ。}\\定義としては、Real-World, Aggregation, Information Model, Decision Model, Implementation の向きを正確に説明する。。具体例では、配送・在庫・フードデリバリーのような現実問題を用い、状態、決定、外生情報、報酬、または情報モデル/意思決定モデルのどこに対応するかを明示する。モデル上の意味として、この概念が目的関数、制約、状態表現、将来価値、または方策選択にどのような影響を与えるかを書く。試験では、用語だけを列挙せず、入力データから意思決定、さらに現実の実装へつながる流れを文章で示すことが重要である。\item \textbf{L01-S016: このスライドの考え方を、講義全体のDDDMプロセスの中に位置づけよ。}\\DDDMプロセスでは、現実世界のデータを集約して情報モデルを作り、意思決定モデルまたは方策で行動を選び、実装後に結果を観測して次の状態へ進む。このスライドの概念は、そのどこか一部だけではなく、データ表現、意思決定、将来不確実性、計算可能性のいずれかに関わる。答案では、まずこの概念がどの段階に属するかを述べ、次に他の段階へどう影響するかを書く。例えば価値関数なら将来価値を現在決定へ戻す役割、FIFOなら旅行時間情報モデルの整合性を保証する役割、CFAなら目的関数を調整して将来に良い行動を誘導する役割である。\end{enumerate}
\end{minipage}
\end{adjustbox}
\end{minipage}


\clearpage
\SlideHeader{Lecture 1 - Overview}{017}{DDDMにおける意思決定}
\begin{minipage}[t]{0.405\textwidth}
\SlideImage{17}{../Materials/2026/3DM_01_Overview_SS26.pdf}
\begin{adjustbox}{max totalsize={\textwidth}{0.43\textheight},center}
\begin{minipage}{\textwidth}\scriptsize
\NoteSection{日本語訳}\begin{itemize}
\item DDDMでは、情報モデルだけを詳しく扱うわけではない。
\item 最適化モデルだけに焦点を当てるわけでもない。
\item 現実世界のプロセスを含むループ全体を考慮する。
\item そのため、システムを時間を通じてモデル化する必要がある。これは dynamic である。
\item 現実世界は外部効果により変化し得る。これは stochastic である。
\item 実装された決定は、すでに変化した現実世界と衝突する可能性がある。
\end{itemize}\NoteSection{試験対策ポイント}\begin{itemize}
\item Real-World, Aggregation, Information Model, Decision Model, Implementation の向きを正確に説明する。
\item 情報モデルは現実の圧縮表現、意思決定モデルは行動選択の評価構造である。
\end{itemize}
\end{minipage}
\end{adjustbox}
\end{minipage}\hfill
\begin{minipage}[t]{0.58\textwidth}
\begin{adjustbox}{max totalsize={\textwidth}{0.89\textheight},center}
\begin{minipage}{\textwidth}\scriptsize
\NoteSection{解説}このページでは「DDDMにおける意思決定」を理解する。スライド上の表現は短いが、試験では定義、具体例、モデル上の意味、手法の限界まで説明できる必要がある。\par
Real-World は、配送車、顧客、道路、倉庫、ドライバー、注文、交通といった現実そのものである。現実は複雑すぎるため、そのまま数理モデルやアルゴリズムに入れることはできない。そこで、意思決定に必要な要素だけを抽出し、モデルが扱える形に変換する必要がある。\par
Aggregation は、現実から得られた詳細データを情報モデルの次元に圧縮する操作である。例として、GPS点列を道路セグメントの速度に変換する、住所を顧客ノードに変換する、過去走行ログから顧客間旅行時間行列を作る、注文履歴から需要分布を推定する、といった処理がある。\par
Information Model は、意思決定に必要な情報の構造化表現である。ここには、顧客集合、デポ、車両集合、旅行時間行列、需要、時間窓、確率分布などが含まれる。重要なのは、情報モデルは現実の全コピーではなく、意思決定に有用な抽象化だという点である。\par
Decision Model は、その情報モデルを入力として、どの行動を選べるか、何を最適化するか、どの制約を守るかを定義する。TSPなら、顧客を一度ずつ訪問し総旅行時間を最小化する。VRPなら、複数車両、容量制約、時間窓などが加わる。ここでは決定変数、目的関数、制約が中心になる。\par
Implementation は、モデル上の解を現実の行動に戻す処理である。たとえば x\_\{ij\}=1 というアーク選択を、ドライバーに提示する訪問順、住所、ナビゲーション、積み込み順へ変換する。実装後には実際の遅延や走行時間が観測され、それが再びAggregationを通じてモデル改善に使われる。\par\NoteSection{確認問題と解答}\begin{enumerate}\item \textbf{L01-S017: 「DDDMにおける意思決定」の概念を、定義・具体例・モデル上の意味の順に説明せよ。}\\定義としては、Real-World, Aggregation, Information Model, Decision Model, Implementation の向きを正確に説明する。。具体例では、配送・在庫・フードデリバリーのような現実問題を用い、状態、決定、外生情報、報酬、または情報モデル/意思決定モデルのどこに対応するかを明示する。モデル上の意味として、この概念が目的関数、制約、状態表現、将来価値、または方策選択にどのような影響を与えるかを書く。試験では、用語だけを列挙せず、入力データから意思決定、さらに現実の実装へつながる流れを文章で示すことが重要である。\item \textbf{L01-S017: このスライドの考え方を、講義全体のDDDMプロセスの中に位置づけよ。}\\DDDMプロセスでは、現実世界のデータを集約して情報モデルを作り、意思決定モデルまたは方策で行動を選び、実装後に結果を観測して次の状態へ進む。このスライドの概念は、そのどこか一部だけではなく、データ表現、意思決定、将来不確実性、計算可能性のいずれかに関わる。答案では、まずこの概念がどの段階に属するかを述べ、次に他の段階へどう影響するかを書く。例えば価値関数なら将来価値を現在決定へ戻す役割、FIFOなら旅行時間情報モデルの整合性を保証する役割、CFAなら目的関数を調整して将来に良い行動を誘導する役割である。\end{enumerate}
\end{minipage}
\end{adjustbox}
\end{minipage}


\clearpage
\SlideHeader{Lecture 1 - Overview}{018}{アジェンダ}
\begin{minipage}[t]{0.405\textwidth}
\SlideImage{18}{../Materials/2026/3DM_01_Overview_SS26.pdf}
\begin{adjustbox}{max totalsize={\textwidth}{0.43\textheight},center}
\begin{minipage}{\textwidth}\scriptsize
\NoteSection{日本語訳}\begin{itemize}
\item 動機
\item Operations Research と Business Analytics の復習
\item モデリングと講義全体の見取り図
\end{itemize}
\end{minipage}
\end{adjustbox}
\end{minipage}\hfill
\begin{minipage}[t]{0.58\textwidth}
\begin{adjustbox}{max totalsize={\textwidth}{0.89\textheight},center}
\begin{minipage}{\textwidth}\scriptsize
\NoteSection{注記}このページは表紙・目次・事務情報・導入画像が中心であるため、無理に確認問題や過去問を付けない。
\end{minipage}
\end{adjustbox}
\end{minipage}


\clearpage
\SlideHeader{Lecture 1 - Overview}{019}{例: Braunschweigの小包配送}
\begin{minipage}[t]{0.405\textwidth}
\SlideImage{19}{../Materials/2026/3DM_01_Overview_SS26.pdf}
\begin{adjustbox}{max totalsize={\textwidth}{0.43\textheight},center}
\begin{minipage}{\textwidth}\scriptsize
\NoteSection{日本語訳}\begin{itemize}
\item トラックはデポにいる。
\item 顧客と住所がある。
\item 道路ネットワークがある。
\item 全顧客を訪問するように、市内でのドライバーのナビゲーションを探す。
\item 目標は迅速な配送である。
\item どうすれば目標を達成できるか。
\end{itemize}\NoteSection{試験対策ポイント}\begin{itemize}
\item TSPの要素をノード、アーク、コスト、決定変数、目的関数、制約に分けて説明する。
\item 時間依存TSPでは d\_ij が d\_ij(t) になり、到着時刻が決定に影響する。
\end{itemize}
\end{minipage}
\end{adjustbox}
\end{minipage}\hfill
\begin{minipage}[t]{0.58\textwidth}
\begin{adjustbox}{max totalsize={\textwidth}{0.89\textheight},center}
\begin{minipage}{\textwidth}\scriptsize
\NoteSection{解説}このページでは「例: Braunschweigの小包配送」を理解する。スライド上の表現は短いが、試験では定義、具体例、モデル上の意味、手法の限界まで説明できる必要がある。\par
Travelling Salesman Problem は、複数の顧客を一度ずつ訪問し、最後に出発点へ戻る最短ツアーを求める基本的なルーティング問題である。顧客やデポはノード、ノード間の移動はアーク、距離や旅行時間はアークコストとして表される。\par
意思決定モデルとしてのTSPでは、x\_\{ij\}=1 なら i から j へ直接移動する、0なら移動しない、という二値変数を使う。目的関数は、選ばれたアークの距離または旅行時間の合計を最小化する。制約として、各顧客に一度入る、各顧客から一度出る、部分巡回を禁止する条件が必要である。\par
情報モデルとしては、顧客集合、デポ、顧客間距離または旅行時間行列が必要である。つまり、TSPは情報モデルと意思決定モデルの典型的な接続例である。旅行時間行列が情報モデル、二値変数・目的関数・制約を持つ最適化問題が意思決定モデルである。\par
時間依存TSPでは、アークコストが出発時刻に依存する。固定の d\_\{ij\} ではなく d\_\{ij\}(t) になるため、同じ i から j への移動でも朝は20分、昼は10分、夕方は30分ということがある。この場合、訪問順序だけでなく、いつそのアークを通るかが重要になる。\par
具体例として、顧客Aを先に訪問するとBへの移動がラッシュ前に終わるが、Cを先に訪問するとBへの移動がラッシュ中になり遅れることがある。つまり、局所的に短い移動を選んでも、時刻の変化により後続コストが悪化するため、時間を含む意思決定モデルが必要になる。\par\NoteSection{確認問題と解答}\begin{enumerate}\item \textbf{L01-S019: 「例: Braunschweigの小包配送」の概念を、定義・具体例・モデル上の意味の順に説明せよ。}\\定義としては、TSPの要素をノード、アーク、コスト、決定変数、目的関数、制約に分けて説明する。。具体例では、配送・在庫・フードデリバリーのような現実問題を用い、状態、決定、外生情報、報酬、または情報モデル/意思決定モデルのどこに対応するかを明示する。モデル上の意味として、この概念が目的関数、制約、状態表現、将来価値、または方策選択にどのような影響を与えるかを書く。試験では、用語だけを列挙せず、入力データから意思決定、さらに現実の実装へつながる流れを文章で示すことが重要である。\item \textbf{L01-S019: このスライドの考え方を、講義全体のDDDMプロセスの中に位置づけよ。}\\DDDMプロセスでは、現実世界のデータを集約して情報モデルを作り、意思決定モデルまたは方策で行動を選び、実装後に結果を観測して次の状態へ進む。このスライドの概念は、そのどこか一部だけではなく、データ表現、意思決定、将来不確実性、計算可能性のいずれかに関わる。答案では、まずこの概念がどの段階に属するかを述べ、次に他の段階へどう影響するかを書く。例えば価値関数なら将来価値を現在決定へ戻す役割、FIFOなら旅行時間情報モデルの整合性を保証する役割、CFAなら目的関数を調整して将来に良い行動を誘導する役割である。\end{enumerate}
\end{minipage}
\end{adjustbox}
\end{minipage}


\clearpage
\SlideHeader{Lecture 1 - Overview}{020}{モデリング: 巡回セールスマン問題}
\begin{minipage}[t]{0.405\textwidth}
\SlideImage{20}{../Materials/2026/3DM_01_Overview_SS26.pdf}
\begin{adjustbox}{max totalsize={\textwidth}{0.43\textheight},center}
\begin{minipage}{\textwidth}\scriptsize
\NoteSection{日本語訳}\begin{itemize}
\item Part 1
\item Part 2
\end{itemize}\NoteSection{試験対策ポイント}\begin{itemize}
\item TSPの要素をノード、アーク、コスト、決定変数、目的関数、制約に分けて説明する。
\item 時間依存TSPでは d\_ij が d\_ij(t) になり、到着時刻が決定に影響する。
\end{itemize}
\end{minipage}
\end{adjustbox}
\end{minipage}\hfill
\begin{minipage}[t]{0.58\textwidth}
\begin{adjustbox}{max totalsize={\textwidth}{0.89\textheight},center}
\begin{minipage}{\textwidth}\scriptsize
\NoteSection{解説}このページでは「モデリング: 巡回セールスマン問題」を理解する。スライド上の表現は短いが、試験では定義、具体例、モデル上の意味、手法の限界まで説明できる必要がある。\par
Travelling Salesman Problem は、複数の顧客を一度ずつ訪問し、最後に出発点へ戻る最短ツアーを求める基本的なルーティング問題である。顧客やデポはノード、ノード間の移動はアーク、距離や旅行時間はアークコストとして表される。\par
意思決定モデルとしてのTSPでは、x\_\{ij\}=1 なら i から j へ直接移動する、0なら移動しない、という二値変数を使う。目的関数は、選ばれたアークの距離または旅行時間の合計を最小化する。制約として、各顧客に一度入る、各顧客から一度出る、部分巡回を禁止する条件が必要である。\par
情報モデルとしては、顧客集合、デポ、顧客間距離または旅行時間行列が必要である。つまり、TSPは情報モデルと意思決定モデルの典型的な接続例である。旅行時間行列が情報モデル、二値変数・目的関数・制約を持つ最適化問題が意思決定モデルである。\par
時間依存TSPでは、アークコストが出発時刻に依存する。固定の d\_\{ij\} ではなく d\_\{ij\}(t) になるため、同じ i から j への移動でも朝は20分、昼は10分、夕方は30分ということがある。この場合、訪問順序だけでなく、いつそのアークを通るかが重要になる。\par
具体例として、顧客Aを先に訪問するとBへの移動がラッシュ前に終わるが、Cを先に訪問するとBへの移動がラッシュ中になり遅れることがある。つまり、局所的に短い移動を選んでも、時刻の変化により後続コストが悪化するため、時間を含む意思決定モデルが必要になる。\par\NoteSection{確認問題と解答}\begin{enumerate}\item \textbf{L01-S020: 「モデリング: 巡回セールスマン問題」の概念を、定義・具体例・モデル上の意味の順に説明せよ。}\\定義としては、TSPの要素をノード、アーク、コスト、決定変数、目的関数、制約に分けて説明する。。具体例では、配送・在庫・フードデリバリーのような現実問題を用い、状態、決定、外生情報、報酬、または情報モデル/意思決定モデルのどこに対応するかを明示する。モデル上の意味として、この概念が目的関数、制約、状態表現、将来価値、または方策選択にどのような影響を与えるかを書く。試験では、用語だけを列挙せず、入力データから意思決定、さらに現実の実装へつながる流れを文章で示すことが重要である。\item \textbf{L01-S020: このスライドの考え方を、講義全体のDDDMプロセスの中に位置づけよ。}\\DDDMプロセスでは、現実世界のデータを集約して情報モデルを作り、意思決定モデルまたは方策で行動を選び、実装後に結果を観測して次の状態へ進む。このスライドの概念は、そのどこか一部だけではなく、データ表現、意思決定、将来不確実性、計算可能性のいずれかに関わる。答案では、まずこの概念がどの段階に属するかを述べ、次に他の段階へどう影響するかを書く。例えば価値関数なら将来価値を現在決定へ戻す役割、FIFOなら旅行時間情報モデルの整合性を保証する役割、CFAなら目的関数を調整して将来に良い行動を誘導する役割である。\end{enumerate}
\end{minipage}
\end{adjustbox}
\end{minipage}


\clearpage
\SlideHeader{Lecture 1 - Overview}{021}{解を見つける}
\begin{minipage}[t]{0.405\textwidth}
\SlideImage{21}{../Materials/2026/3DM_01_Overview_SS26.pdf}
\begin{adjustbox}{max totalsize={\textwidth}{0.43\textheight},center}
\begin{minipage}{\textwidth}\scriptsize
\NoteSection{日本語訳}\begin{itemize}
\item モデルは解空間を張り、解を評価する。
\item 候補解の求め方
\item 最後のステップ: 実務への実装
\end{itemize}\NoteSection{試験対策ポイント}\begin{itemize}
\item TSPの要素をノード、アーク、コスト、決定変数、目的関数、制約に分けて説明する。
\item 時間依存TSPでは d\_ij が d\_ij(t) になり、到着時刻が決定に影響する。
\end{itemize}
\end{minipage}
\end{adjustbox}
\end{minipage}\hfill
\begin{minipage}[t]{0.58\textwidth}
\begin{adjustbox}{max totalsize={\textwidth}{0.89\textheight},center}
\begin{minipage}{\textwidth}\scriptsize
\NoteSection{解説}このページでは「解を見つける」を理解する。スライド上の表現は短いが、試験では定義、具体例、モデル上の意味、手法の限界まで説明できる必要がある。\par
Travelling Salesman Problem は、複数の顧客を一度ずつ訪問し、最後に出発点へ戻る最短ツアーを求める基本的なルーティング問題である。顧客やデポはノード、ノード間の移動はアーク、距離や旅行時間はアークコストとして表される。\par
意思決定モデルとしてのTSPでは、x\_\{ij\}=1 なら i から j へ直接移動する、0なら移動しない、という二値変数を使う。目的関数は、選ばれたアークの距離または旅行時間の合計を最小化する。制約として、各顧客に一度入る、各顧客から一度出る、部分巡回を禁止する条件が必要である。\par
情報モデルとしては、顧客集合、デポ、顧客間距離または旅行時間行列が必要である。つまり、TSPは情報モデルと意思決定モデルの典型的な接続例である。旅行時間行列が情報モデル、二値変数・目的関数・制約を持つ最適化問題が意思決定モデルである。\par
時間依存TSPでは、アークコストが出発時刻に依存する。固定の d\_\{ij\} ではなく d\_\{ij\}(t) になるため、同じ i から j への移動でも朝は20分、昼は10分、夕方は30分ということがある。この場合、訪問順序だけでなく、いつそのアークを通るかが重要になる。\par
具体例として、顧客Aを先に訪問するとBへの移動がラッシュ前に終わるが、Cを先に訪問するとBへの移動がラッシュ中になり遅れることがある。つまり、局所的に短い移動を選んでも、時刻の変化により後続コストが悪化するため、時間を含む意思決定モデルが必要になる。\par\NoteSection{確認問題と解答}\begin{enumerate}\item \textbf{L01-S021: 「解を見つける」の概念を、定義・具体例・モデル上の意味の順に説明せよ。}\\定義としては、TSPの要素をノード、アーク、コスト、決定変数、目的関数、制約に分けて説明する。。具体例では、配送・在庫・フードデリバリーのような現実問題を用い、状態、決定、外生情報、報酬、または情報モデル/意思決定モデルのどこに対応するかを明示する。モデル上の意味として、この概念が目的関数、制約、状態表現、将来価値、または方策選択にどのような影響を与えるかを書く。試験では、用語だけを列挙せず、入力データから意思決定、さらに現実の実装へつながる流れを文章で示すことが重要である。\item \textbf{L01-S021: このスライドの考え方を、講義全体のDDDMプロセスの中に位置づけよ。}\\DDDMプロセスでは、現実世界のデータを集約して情報モデルを作り、意思決定モデルまたは方策で行動を選び、実装後に結果を観測して次の状態へ進む。このスライドの概念は、そのどこか一部だけではなく、データ表現、意思決定、将来不確実性、計算可能性のいずれかに関わる。答案では、まずこの概念がどの段階に属するかを述べ、次に他の段階へどう影響するかを書く。例えば価値関数なら将来価値を現在決定へ戻す役割、FIFOなら旅行時間情報モデルの整合性を保証する役割、CFAなら目的関数を調整して将来に良い行動を誘導する役割である。\end{enumerate}
\end{minipage}
\end{adjustbox}
\end{minipage}


\clearpage
\SlideHeader{Lecture 1 - Overview}{022}{Operations Researchの三段階}
\begin{minipage}[t]{0.405\textwidth}
\SlideImage{22}{../Materials/2026/3DM_01_Overview_SS26.pdf}
\begin{adjustbox}{max totalsize={\textwidth}{0.43\textheight},center}
\begin{minipage}{\textwidth}\scriptsize
\NoteSection{日本語訳}\begin{itemize}
\item 問題定義
\item モデリング
\item 解法
\item 情報と意思決定
\item 実装を含む。
\end{itemize}\NoteSection{試験対策ポイント}\begin{itemize}
\item IDAは appearance から structure、PMT/ORは structure から solution へ向かう。
\item Business Analytics は Statistics, OR, CS/IS の統合領域である。
\end{itemize}
\end{minipage}
\end{adjustbox}
\end{minipage}\hfill
\begin{minipage}[t]{0.58\textwidth}
\begin{adjustbox}{max totalsize={\textwidth}{0.89\textheight},center}
\begin{minipage}{\textwidth}\scriptsize
\NoteSection{解説}このページでは「Operations Researchの三段階」を理解する。スライド上の表現は短いが、試験では定義、具体例、モデル上の意味、手法の限界まで説明できる必要がある。\par
IDA, Intelligent Data Analysis は、記録されたデータの appearance から出発する。appearance とは、現実システムの表面的な観測結果であり、GPSログ、注文履歴、クリック履歴、センサー値のようなものである。これらは誤差、欠損、外れ値を含むため、そのまま構造を表しているとは限らない。\par
IDAの役割は、データをクリーニングし、補完し、検証し、統計的パターンやモデルを抽出することで、背後にある structure への証拠を与えることである。例えば、タクシーの走行ログから朝夕ラッシュの速度低下パターンを見つけることは、道路交通の構造を推定する作業である。\par
PMTやOperations Researchは逆方向で考える。まず現実問題の重要構造を仮定し、目的、制約、決定空間を明確にする。そして、ルーティング、スケジューリング、割当、在庫管理などの意思決定モデルを作り、解を求める。これは structure から実行可能な appearance、つまり具体的な解へ向かう流れである。\par
Business Analytics の三つの基礎は、Statistics、Computer Science/Information Systems、Operations Research である。Statistics は推定と不確実性、CS/IS はデータ処理と実装、OR は最適化と意思決定を担う。三つを列挙するだけでは不十分で、互いの交差領域が現実のデータ駆動意思決定で不可欠であることを説明する必要がある。\par
具体例として配送計画を考えると、Statistics は旅行時間分布を推定し、CS/IS はGPSデータや注文データを保存・処理し、OR は車両ルートを最適化する。データだけではルートは決まらず、最適化だけでは現実の交通を反映できないため、この統合がDDDMの基礎になる。\par\NoteSection{確認問題と解答}\begin{enumerate}\item \textbf{L01-S022: 「Operations Researchの三段階」の概念を、定義・具体例・モデル上の意味の順に説明せよ。}\\定義としては、IDAは appearance から structure、PMT/ORは structure から solution へ向かう。。具体例では、配送・在庫・フードデリバリーのような現実問題を用い、状態、決定、外生情報、報酬、または情報モデル/意思決定モデルのどこに対応するかを明示する。モデル上の意味として、この概念が目的関数、制約、状態表現、将来価値、または方策選択にどのような影響を与えるかを書く。試験では、用語だけを列挙せず、入力データから意思決定、さらに現実の実装へつながる流れを文章で示すことが重要である。\item \textbf{L01-S022: このスライドの考え方を、講義全体のDDDMプロセスの中に位置づけよ。}\\DDDMプロセスでは、現実世界のデータを集約して情報モデルを作り、意思決定モデルまたは方策で行動を選び、実装後に結果を観測して次の状態へ進む。このスライドの概念は、そのどこか一部だけではなく、データ表現、意思決定、将来不確実性、計算可能性のいずれかに関わる。答案では、まずこの概念がどの段階に属するかを述べ、次に他の段階へどう影響するかを書く。例えば価値関数なら将来価値を現在決定へ戻す役割、FIFOなら旅行時間情報モデルの整合性を保証する役割、CFAなら目的関数を調整して将来に良い行動を誘導する役割である。\end{enumerate}
\end{minipage}
\end{adjustbox}
\end{minipage}


\clearpage
\SlideHeader{Lecture 1 - Overview}{023}{アジェンダ}
\begin{minipage}[t]{0.405\textwidth}
\SlideImage{23}{../Materials/2026/3DM_01_Overview_SS26.pdf}
\begin{adjustbox}{max totalsize={\textwidth}{0.43\textheight},center}
\begin{minipage}{\textwidth}\scriptsize
\NoteSection{日本語訳}\begin{itemize}
\item 動機
\item Operations Research と Business Analytics の復習
\item モデリングと講義全体の見取り図
\end{itemize}
\end{minipage}
\end{adjustbox}
\end{minipage}\hfill
\begin{minipage}[t]{0.58\textwidth}
\begin{adjustbox}{max totalsize={\textwidth}{0.89\textheight},center}
\begin{minipage}{\textwidth}\scriptsize
\NoteSection{注記}このページは表紙・目次・事務情報・導入画像が中心であるため、無理に確認問題や過去問を付けない。
\end{minipage}
\end{adjustbox}
\end{minipage}


\clearpage
\SlideHeader{Lecture 1 - Overview}{024}{モデリング}
\begin{minipage}[t]{0.405\textwidth}
\SlideImage{24}{../Materials/2026/3DM_01_Overview_SS26.pdf}
\begin{adjustbox}{max totalsize={\textwidth}{0.43\textheight},center}
\begin{minipage}{\textwidth}\scriptsize
\NoteSection{日本語訳}\begin{itemize}
\item 異なる問題が同じモデルを共有する場合がある。例: 小包配送と技術者サービス。
\item 同じ問題を異なる方法でモデル化することもできる。
\item 情報モデルと意思決定モデルの間には接続がある。
\end{itemize}\NoteSection{試験対策ポイント}\begin{itemize}
\item Real-World, Aggregation, Information Model, Decision Model, Implementation の向きを正確に説明する。
\item 情報モデルは現実の圧縮表現、意思決定モデルは行動選択の評価構造である。
\end{itemize}
\end{minipage}
\end{adjustbox}
\end{minipage}\hfill
\begin{minipage}[t]{0.58\textwidth}
\begin{adjustbox}{max totalsize={\textwidth}{0.89\textheight},center}
\begin{minipage}{\textwidth}\scriptsize
\NoteSection{解説}このページでは「モデリング」を理解する。スライド上の表現は短いが、試験では定義、具体例、モデル上の意味、手法の限界まで説明できる必要がある。\par
Real-World は、配送車、顧客、道路、倉庫、ドライバー、注文、交通といった現実そのものである。現実は複雑すぎるため、そのまま数理モデルやアルゴリズムに入れることはできない。そこで、意思決定に必要な要素だけを抽出し、モデルが扱える形に変換する必要がある。\par
Aggregation は、現実から得られた詳細データを情報モデルの次元に圧縮する操作である。例として、GPS点列を道路セグメントの速度に変換する、住所を顧客ノードに変換する、過去走行ログから顧客間旅行時間行列を作る、注文履歴から需要分布を推定する、といった処理がある。\par
Information Model は、意思決定に必要な情報の構造化表現である。ここには、顧客集合、デポ、車両集合、旅行時間行列、需要、時間窓、確率分布などが含まれる。重要なのは、情報モデルは現実の全コピーではなく、意思決定に有用な抽象化だという点である。\par
Decision Model は、その情報モデルを入力として、どの行動を選べるか、何を最適化するか、どの制約を守るかを定義する。TSPなら、顧客を一度ずつ訪問し総旅行時間を最小化する。VRPなら、複数車両、容量制約、時間窓などが加わる。ここでは決定変数、目的関数、制約が中心になる。\par
Implementation は、モデル上の解を現実の行動に戻す処理である。たとえば x\_\{ij\}=1 というアーク選択を、ドライバーに提示する訪問順、住所、ナビゲーション、積み込み順へ変換する。実装後には実際の遅延や走行時間が観測され、それが再びAggregationを通じてモデル改善に使われる。\par\NoteSection{確認問題と解答}\begin{enumerate}\item \textbf{L01-S024: 「モデリング」の概念を、定義・具体例・モデル上の意味の順に説明せよ。}\\定義としては、Real-World, Aggregation, Information Model, Decision Model, Implementation の向きを正確に説明する。。具体例では、配送・在庫・フードデリバリーのような現実問題を用い、状態、決定、外生情報、報酬、または情報モデル/意思決定モデルのどこに対応するかを明示する。モデル上の意味として、この概念が目的関数、制約、状態表現、将来価値、または方策選択にどのような影響を与えるかを書く。試験では、用語だけを列挙せず、入力データから意思決定、さらに現実の実装へつながる流れを文章で示すことが重要である。\item \textbf{L01-S024: このスライドの考え方を、講義全体のDDDMプロセスの中に位置づけよ。}\\DDDMプロセスでは、現実世界のデータを集約して情報モデルを作り、意思決定モデルまたは方策で行動を選び、実装後に結果を観測して次の状態へ進む。このスライドの概念は、そのどこか一部だけではなく、データ表現、意思決定、将来不確実性、計算可能性のいずれかに関わる。答案では、まずこの概念がどの段階に属するかを述べ、次に他の段階へどう影響するかを書く。例えば価値関数なら将来価値を現在決定へ戻す役割、FIFOなら旅行時間情報モデルの整合性を保証する役割、CFAなら目的関数を調整して将来に良い行動を誘導する役割である。\end{enumerate}
\end{minipage}
\end{adjustbox}
\end{minipage}


\clearpage
\SlideHeader{Lecture 1 - Overview}{025}{例I: 代替的な情報モデル}
\begin{minipage}[t]{0.405\textwidth}
\SlideImage{25}{../Materials/2026/3DM_01_Overview_SS26.pdf}
\begin{adjustbox}{max totalsize={\textwidth}{0.43\textheight},center}
\begin{minipage}{\textwidth}\scriptsize
\NoteSection{日本語訳}\begin{itemize}
\item 旅行時間の日内変動、つまり時間依存旅行時間を観測する。
\item 7-9時、9-15時、15-18時のような時間帯を考える。
\item 情報モデルは複数の旅行時間行列の集合を含む。
\item 情報生成はより複雑になる。
\item 代替的な意思決定モデルが必要になる。
\end{itemize}\NoteSection{試験対策ポイント}\begin{itemize}
\item 時間依存旅行時間は情報モデルを精密にするが、意思決定モデルも複雑にする。
\item 平均行列、時間帯別行列、連続関数の違いを説明できる。
\end{itemize}
\end{minipage}
\end{adjustbox}
\end{minipage}\hfill
\begin{minipage}[t]{0.58\textwidth}
\begin{adjustbox}{max totalsize={\textwidth}{0.89\textheight},center}
\begin{minipage}{\textwidth}\scriptsize
\NoteSection{解説}このページでは「例I: 代替的な情報モデル」を理解する。スライド上の表現は短いが、試験では定義、具体例、モデル上の意味、手法の限界まで説明できる必要がある。\par
時間依存旅行時間とは、同じ道路や同じ顧客間でも、出発時刻によって旅行時間が変わることを意味する。通勤時間帯、昼間、夜間、週末で交通量が変わるため、単一の平均旅行時間では現実を十分に表せない場合が多い。\par
最も単純な情報モデルは、すべての時刻に共通する平均旅行時間行列である。これは扱いやすいが、朝夕の混雑や夜間の自由流交通を無視する。次に、時間帯別旅行時間行列がある。これは、例えば月曜8時台、月曜9時台のように複数の行列を用意する方法である。\par
さらに詳細には、各道路セグメントや顧客間アークに対して連続的な旅行時間関数を持つこともできる。ただし、細かいモデルほど多くのデータが必要になり、各時間帯・道路ごとの観測数が少なくなると推定が不安定になる。詳細さは常に良いわけではなく、意思決定に役立つ粒度を選ぶ必要がある。\par
時間依存性は意思決定モデルにも影響する。静的TSPではアークコストは固定なので訪問順だけを考えればよいが、時間依存TSPでは到着時刻、出発時刻、サービス時間、次アークの旅行時間が連鎖する。したがって、ある顧客を挿入するだけで後続すべての到着時刻が変わる。\par
具体例として、昼に郊外を先に回り、夕方に市中心部を回る計画と、逆の計画では、同じ顧客を訪問していても総旅行時間が大きく異なる。試験答案では、時間依存性が単なるデータの違いではなく、ルーティング決定の構造を変える点まで述べるとよい。\par\NoteSection{確認問題と解答}\begin{enumerate}\item \textbf{L01-S025: 「例I: 代替的な情報モデル」の概念を、定義・具体例・モデル上の意味の順に説明せよ。}\\定義としては、時間依存旅行時間は情報モデルを精密にするが、意思決定モデルも複雑にする。。具体例では、配送・在庫・フードデリバリーのような現実問題を用い、状態、決定、外生情報、報酬、または情報モデル/意思決定モデルのどこに対応するかを明示する。モデル上の意味として、この概念が目的関数、制約、状態表現、将来価値、または方策選択にどのような影響を与えるかを書く。試験では、用語だけを列挙せず、入力データから意思決定、さらに現実の実装へつながる流れを文章で示すことが重要である。\item \textbf{L01-S025: このスライドの考え方を、講義全体のDDDMプロセスの中に位置づけよ。}\\DDDMプロセスでは、現実世界のデータを集約して情報モデルを作り、意思決定モデルまたは方策で行動を選び、実装後に結果を観測して次の状態へ進む。このスライドの概念は、そのどこか一部だけではなく、データ表現、意思決定、将来不確実性、計算可能性のいずれかに関わる。答案では、まずこの概念がどの段階に属するかを述べ、次に他の段階へどう影響するかを書く。例えば価値関数なら将来価値を現在決定へ戻す役割、FIFOなら旅行時間情報モデルの整合性を保証する役割、CFAなら目的関数を調整して将来に良い行動を誘導する役割である。\end{enumerate}
\end{minipage}
\end{adjustbox}
\end{minipage}


\clearpage
\SlideHeader{Lecture 1 - Overview}{026}{例II: 代替的な意思決定モデル}
\begin{minipage}[t]{0.405\textwidth}
\SlideImage{26}{../Materials/2026/3DM_01_Overview_SS26.pdf}
\begin{adjustbox}{max totalsize={\textwidth}{0.43\textheight},center}
\begin{minipage}{\textwidth}\scriptsize
\NoteSection{日本語訳}\begin{itemize}
\item 解が実行可能に見えても、残業が頻繁に発生する。
\item 旅行時間に不確実性がある。
\item 意思決定モデルを変更する。これは stochasticity に関係する。
\item 代替的でより詳細な情報モデルが必要になる。
\end{itemize}\NoteSection{試験対策ポイント}\begin{itemize}
\item 不確実性は、需要・リソース・環境の変化として整理できる。
\item 頻度、範囲、破壊力、予測可能性で性質が異なる。
\end{itemize}
\end{minipage}
\end{adjustbox}
\end{minipage}\hfill
\begin{minipage}[t]{0.58\textwidth}
\begin{adjustbox}{max totalsize={\textwidth}{0.89\textheight},center}
\begin{minipage}{\textwidth}\scriptsize
\NoteSection{解説}このページでは「例II: 代替的な意思決定モデル」を理解する。スライド上の表現は短いが、試験では定義、具体例、モデル上の意味、手法の限界まで説明できる必要がある。\par
不確実性とは、意思決定時点では将来の状態や入力値が完全には分からないことである。配送であれば、注文がどこから来るか、交通がどれだけ混むか、サービス時間が何分か、ドライバーが利用可能かどうかが事前には確定していない。\par
不確実性は一種類ではない。需要不確実性は、顧客からの注文の発生時刻、場所、量に関する不確実性である。リソース不確実性は、車両故障、ドライバーのキャンセル、バッテリー残量、積載能力に関する不確実性である。環境不確実性は、交通、天候、道路工事、事故、駐車可能性に関する不確実性である。\par
不確実性は、頻度、範囲、破壊力、予測可能性によって性質が異なる。例えば日々の交通変動は頻繁に起きるが、過去データからある程度予測できる。一方で、大規模事故や道路閉鎖は低頻度だが、発生すれば計画を大きく壊す。\par
不確実性を無視して平均値だけで計画すると、平均的には良く見えるが実行時に頻繁に遅れる計画になる可能性がある。特に複数顧客を連続して訪問するルーティングでは、一つの遅延が後続訪問に波及するため、単独の平均時間以上にリスクが大きい。\par
試験答案では、不確実性の例を挙げるだけでなく、それが情報モデルと意思決定モデルにどう影響するかを書くとよい。情報モデルでは確率分布や区間として表現し、意思決定モデルでは期待値、リスク、ロバスト性、再最適化などを考える。\par\NoteSection{確認問題と解答}\begin{enumerate}\item \textbf{L01-S026: 「例II: 代替的な意思決定モデル」の概念を、定義・具体例・モデル上の意味の順に説明せよ。}\\定義としては、不確実性は、需要・リソース・環境の変化として整理できる。。具体例では、配送・在庫・フードデリバリーのような現実問題を用い、状態、決定、外生情報、報酬、または情報モデル/意思決定モデルのどこに対応するかを明示する。モデル上の意味として、この概念が目的関数、制約、状態表現、将来価値、または方策選択にどのような影響を与えるかを書く。試験では、用語だけを列挙せず、入力データから意思決定、さらに現実の実装へつながる流れを文章で示すことが重要である。\item \textbf{L01-S026: このスライドの考え方を、講義全体のDDDMプロセスの中に位置づけよ。}\\DDDMプロセスでは、現実世界のデータを集約して情報モデルを作り、意思決定モデルまたは方策で行動を選び、実装後に結果を観測して次の状態へ進む。このスライドの概念は、そのどこか一部だけではなく、データ表現、意思決定、将来不確実性、計算可能性のいずれかに関わる。答案では、まずこの概念がどの段階に属するかを述べ、次に他の段階へどう影響するかを書く。例えば価値関数なら将来価値を現在決定へ戻す役割、FIFOなら旅行時間情報モデルの整合性を保証する役割、CFAなら目的関数を調整して将来に良い行動を誘導する役割である。\end{enumerate}
\end{minipage}
\end{adjustbox}
\end{minipage}


\clearpage
\SlideHeader{Lecture 1 - Overview}{027}{例III: 動的性}
\begin{minipage}[t]{0.405\textwidth}
\SlideImage{27}{../Materials/2026/3DM_01_Overview_SS26.pdf}
\begin{adjustbox}{max totalsize={\textwidth}{0.43\textheight},center}
\begin{minipage}{\textwidth}\scriptsize
\NoteSection{日本語訳}\begin{itemize}
\item 情報モデルが時間とともに変化することを観測する。
\item 変化に反応できる。
\item 確率的・動的モデルが必要になる。
\item 将来の先読みが重要になる。
\item Approximate Dynamic Programming へつながる。
\end{itemize}\NoteSection{試験対策ポイント}\begin{itemize}
\item DDDMはデータ分析を意思決定と実装へ接続する講義である。
\item stochasticity と dynamism を分けて説明できることが最重要である。
\end{itemize}
\end{minipage}
\end{adjustbox}
\end{minipage}\hfill
\begin{minipage}[t]{0.58\textwidth}
\begin{adjustbox}{max totalsize={\textwidth}{0.89\textheight},center}
\begin{minipage}{\textwidth}\scriptsize
\NoteSection{解説}このページでは「例III: 動的性」を理解する。スライド上の表現は短いが、試験では定義、具体例、モデル上の意味、手法の限界まで説明できる必要がある。\par
Data Driven Decision Making は、単にデータを分析してレポートを作ることではない。データを使って、現実で実行する行動を選び、その結果をまた観測して次の意思決定へ反映する考え方である。たとえば食品配送であれば、注文数を予測するだけではなく、どの注文を受けるか、どのドライバーに割り当てるか、どの順番で回るかまで決める。\par
この講義の中心は、現実の意思決定が静的で完全情報の問題ではないという点にある。現実では、注文、交通、ドライバー位置、キャンセル、サービス時間などが時間とともに変わる。したがって、最初に作った計画を最後まで固定するのではなく、新しい情報を受けて状態を更新し、必要に応じて計画を変える必要がある。\par
stochasticity は、将来何が起こるかが確率的・不確実であることを意味する。dynamism は、時間が進むにつれて状態が変わり、意思決定を繰り返す必要があることを意味する。試験では、この二語を同義語のように書かないことが重要である。注文がいつ来るか分からないことは stochasticity、注文が来た後に計画を更新することは dynamism である。\par
答案では、DDDMを Business Analytics、Operations Research、Data Analytics、情報システム実装の交点として説明すると強い。予測モデルだけでは『何をするか』が決まらず、最適化モデルだけでは『現実データをどう入れるか』が決まらない。DDDMはこの間を接続し、現実のデータから実行可能な意思決定を作る。\par
具体例として、Uberや食品配送では、リアルタイム注文、ドライバーの現在位置、交通状況、顧客の待ち時間を使いながら、数分ごとに新しい割当を作る。このとき、過去データから需要を学習するだけでなく、将来の不確実な注文を見越して現在のドライバーを使い切らないようにするなど、将来の柔軟性も重要になる。\par\NoteSection{確認問題と解答}\begin{enumerate}\item \textbf{L01-S027: 「例III: 動的性」の概念を、定義・具体例・モデル上の意味の順に説明せよ。}\\定義としては、DDDMはデータ分析を意思決定と実装へ接続する講義である。。具体例では、配送・在庫・フードデリバリーのような現実問題を用い、状態、決定、外生情報、報酬、または情報モデル/意思決定モデルのどこに対応するかを明示する。モデル上の意味として、この概念が目的関数、制約、状態表現、将来価値、または方策選択にどのような影響を与えるかを書く。試験では、用語だけを列挙せず、入力データから意思決定、さらに現実の実装へつながる流れを文章で示すことが重要である。\item \textbf{L01-S027: このスライドの考え方を、講義全体のDDDMプロセスの中に位置づけよ。}\\DDDMプロセスでは、現実世界のデータを集約して情報モデルを作り、意思決定モデルまたは方策で行動を選び、実装後に結果を観測して次の状態へ進む。このスライドの概念は、そのどこか一部だけではなく、データ表現、意思決定、将来不確実性、計算可能性のいずれかに関わる。答案では、まずこの概念がどの段階に属するかを述べ、次に他の段階へどう影響するかを書く。例えば価値関数なら将来価値を現在決定へ戻す役割、FIFOなら旅行時間情報モデルの整合性を保証する役割、CFAなら目的関数を調整して将来に良い行動を誘導する役割である。\end{enumerate}
\end{minipage}
\end{adjustbox}
\end{minipage}


\clearpage
\SlideHeader{Lecture 1 - Overview}{028}{次回予告}
\begin{minipage}[t]{0.405\textwidth}
\SlideImage{28}{../Materials/2026/3DM_01_Overview_SS26.pdf}
\begin{adjustbox}{max totalsize={\textwidth}{0.43\textheight},center}
\begin{minipage}{\textwidth}\scriptsize
\NoteSection{日本語訳}\begin{itemize}
\item Business Analytics
\item 情報モデリングと意思決定モデリング
\item ケーススタディ: 時間依存旅行時間を持つ配送ルーティング
\end{itemize}
\end{minipage}
\end{adjustbox}
\end{minipage}\hfill
\begin{minipage}[t]{0.58\textwidth}
\begin{adjustbox}{max totalsize={\textwidth}{0.89\textheight},center}
\begin{minipage}{\textwidth}\scriptsize
\NoteSection{注記}このページは表紙・目次・事務情報・導入画像が中心であるため、無理に確認問題や過去問を付けない。
\end{minipage}
\end{adjustbox}
\end{minipage}

\clearpage\ExamHeader{Lecture 1 - Overview}
\ExamQuestion{SS23 Aufgabe 5 (8 点)}
\ExamSubsection{問題内容}
レストラン配送のCFA/PFA。設問では、遅延リスクを費用近似または方策近似で扱い、注文統合の判断を説明すること。 ことが求められる。\par
\ExamSubsection{満点答案}
満点答案では、Real Worldから観測データを集約してInformation Modelを作り、そこからDecision Modelを作り、解をImplementation/Disaggregationで現実の行動へ戻す流れを説明する。Information Modelは現実の圧縮表現であり、Decision Modelは決定変数、目的関数、制約により行動選択を評価する構造である。\par
\ExamSubsection{具体例による解説}
具体例として、配送問題なら、顧客注文・車両位置・旅行時間を情報モデルにし、訪問順や割当を意思決定モデルで決める。在庫問題なら、在庫水準を状態、注文量を決定、需要を外生情報、在庫更新式を遷移、欠品費・保管費を費用として整理する。問題文の文脈に合わせて、抽象用語を必ず具体的な変数や行動に置き換える。\par
\ExamSubsection{採点で落とさない要素}
\begin{itemize}
\item 設問が求める概念の定義を最初に明確に書く。
\item 講義の専門語を列挙するだけでなく、現実例とモデル要素へ対応付ける。
\item 利点だけでなく限界・注意点も最低一つ述べる。
\item 式が出る問題では、記号の意味を文章で説明する。
\item 新しい事例問題では、状態・決定・外生情報・遷移・報酬/費用の順に整理する。
\end{itemize}
\vspace{2mm}\hrule\vspace{2mm}
\ExamQuestion{WS23/24 Aufgabe 2 (5 点)}
\ExamSubsection{問題内容}
Real-World ProblemとDecision Modelの関係。設問では、Real World, Aggregation, Information Model, Decision Model, Implementation/Disaggregationのデータフローを説明すること。 ことが求められる。\par
\ExamSubsection{満点答案}
満点答案では、Real Worldから観測データを集約してInformation Modelを作り、そこからDecision Modelを作り、解をImplementation/Disaggregationで現実の行動へ戻す流れを説明する。Information Modelは現実の圧縮表現であり、Decision Modelは決定変数、目的関数、制約により行動選択を評価する構造である。\par
\ExamSubsection{具体例による解説}
具体例として、配送問題なら、顧客注文・車両位置・旅行時間を情報モデルにし、訪問順や割当を意思決定モデルで決める。在庫問題なら、在庫水準を状態、注文量を決定、需要を外生情報、在庫更新式を遷移、欠品費・保管費を費用として整理する。問題文の文脈に合わせて、抽象用語を必ず具体的な変数や行動に置き換える。\par
\ExamSubsection{採点で落とさない要素}
\begin{itemize}
\item 設問が求める概念の定義を最初に明確に書く。
\item 講義の専門語を列挙するだけでなく、現実例とモデル要素へ対応付ける。
\item 利点だけでなく限界・注意点も最低一つ述べる。
\item 式が出る問題では、記号の意味を文章で説明する。
\item 新しい事例問題では、状態・決定・外生情報・遷移・報酬/費用の順に整理する。
\end{itemize}
\vspace{2mm}\hrule\vspace{2mm}
\ExamQuestion{SS24 Aufgabe 1a (8 点)}
\ExamSubsection{問題内容}
Information Model, Decision Model, Aggregation, Disaggregationの機能。設問では、各ブロックの役割を定義し、データがどの向きに流れるかを説明すること。 ことが求められる。\par
\ExamSubsection{満点答案}
満点答案では、Real Worldから観測データを集約してInformation Modelを作り、そこからDecision Modelを作り、解をImplementation/Disaggregationで現実の行動へ戻す流れを説明する。Information Modelは現実の圧縮表現であり、Decision Modelは決定変数、目的関数、制約により行動選択を評価する構造である。\par
\ExamSubsection{具体例による解説}
具体例として、配送問題なら、顧客注文・車両位置・旅行時間を情報モデルにし、訪問順や割当を意思決定モデルで決める。在庫問題なら、在庫水準を状態、注文量を決定、需要を外生情報、在庫更新式を遷移、欠品費・保管費を費用として整理する。問題文の文脈に合わせて、抽象用語を必ず具体的な変数や行動に置き換える。\par
\ExamSubsection{採点で落とさない要素}
\begin{itemize}
\item 設問が求める概念の定義を最初に明確に書く。
\item 講義の専門語を列挙するだけでなく、現実例とモデル要素へ対応付ける。
\item 利点だけでなく限界・注意点も最低一つ述べる。
\item 式が出る問題では、記号の意味を文章で説明する。
\item 新しい事例問題では、状態・決定・外生情報・遷移・報酬/費用の順に整理する。
\end{itemize}
\vspace{2mm}\hrule\vspace{2mm}
\ExamQuestion{SS25 Exercise 10 (12 点)}
\ExamSubsection{問題内容}
ADP手法の2つ組み合わせを3例提案。設問では、組み合わせの狙い、補完関係、期待利点を説明すること。 ことが求められる。\par
\ExamSubsection{満点答案}
満点答案では、PFA、Rolling Horizon、CFA、Look-ahead、VFAを列挙し、それぞれが何を近似するかを説明する。PFAは方策そのもの、CFAは目的関数や制約、Look-aheadは将来シナリオ評価、VFAは将来価値を近似する。各手法の強みと弱みも必ず書く。\par
\ExamSubsection{具体例による解説}
具体例として、配送問題なら、顧客注文・車両位置・旅行時間を情報モデルにし、訪問順や割当を意思決定モデルで決める。在庫問題なら、在庫水準を状態、注文量を決定、需要を外生情報、在庫更新式を遷移、欠品費・保管費を費用として整理する。問題文の文脈に合わせて、抽象用語を必ず具体的な変数や行動に置き換える。\par
\ExamSubsection{採点で落とさない要素}
\begin{itemize}
\item 設問が求める概念の定義を最初に明確に書く。
\item 講義の専門語を列挙するだけでなく、現実例とモデル要素へ対応付ける。
\item 利点だけでなく限界・注意点も最低一つ述べる。
\item 式が出る問題では、記号の意味を文章で説明する。
\item 新しい事例問題では、状態・決定・外生情報・遷移・報酬/費用の順に整理する。
\end{itemize}
\vspace{2mm}\hrule\vspace{2mm}